\section{Hybrid MPI + Local Monitoring Approach}
\label{sec:hybrid}

\subsection{The PCAP Capture Challenge}

A critical but often overlooked aspect of network simulation is \textit{observability}—the ability to inspect and analyze protocol behavior. Researchers rely heavily on packet capture (PCAP) files to:

\begin{itemize}
    \item Validate protocol implementations against standards
    \item Debug unexpected behavior through packet-level inspection
    \item Analyze traffic patterns and timing characteristics
    \item Generate publication-quality protocol traces
    \item Demonstrate correctness to reviewers and collaborators
\end{itemize}

NS-3's WiFi module provides excellent PCAP support through the \texttt{YansWifiPhyHelper::EnablePcap()} method, which captures all frames at the PHY layer including 802.11 management frames, control frames, and data frames with complete headers.

However, our MPI-based approach creates a fundamental problem: \textbf{MPI stubs don't generate local WiFi frames}. When \texttt{RemoteYansWifiChannelStub::Send()} is called, instead of transmitting a packet through a local WiFi channel (which PCAP could capture), it serializes the packet into an MPI message and sends it to rank 0. The packet never appears on a local network interface—it's transmitted as an opaque byte buffer between processes. Standard PCAP mechanisms capture nothing.

This creates an unacceptable situation: distributed simulations gain performance but lose observability. Researchers must choose between scalability and debuggability.

\subsection{Dual-Network Solution}

We solve this problem through a \textit{hybrid architecture} where each device simultaneously participates in two independent networks, as shown in Figure~\ref{fig:hybrid-architecture}:

\begin{figure}[htbp]
\centering
\begin{verbatim}
┌─────────────────────────────────────────────────────────────────┐
│                    Device Node (e.g., Rank 1)                   │
│                                                                  │
│  ┌────────────────────────────────────────────────────────────┐ │
│  │                  Application Layer                          │ │
│  │         (UDP Echo Client/Server, OnOff, etc.)              │ │
│  └──────────────────────┬─────────────────┬───────────────────┘ │
│                         │                 │                      │
│            ┌────────────▼─────┐    ┌─────▼──────────┐          │
│            │ IP Stack 1       │    │ IP Stack 2     │          │
│            │ 192.168.x.x      │    │ 10.x.x.x       │          │
│            └────────┬─────────┘    └──────┬─────────┘          │
│                     │                      │                     │
│      ┌──────────────▼──────────┐  ┌───────▼──────────────┐     │
│      │  WiFi Device 0          │  │  WiFi Device 1        │     │
│      │  (MPI Performance)      │  │  (Local Monitoring)   │     │
│      ├─────────────────────────┤  ├───────────────────────┤     │
│      │  MAC: AdHoc Mode        │  │  MAC: AdHoc Mode      │     │
│      │  PHY: High TX Power     │  │  PHY: Low TX Power    │     │
│      └──────────┬──────────────┘  └──────┬────────────────┘     │
│                 │                         │                      │
│      ┌──────────▼────────────────┐  ┌────▼─────────────────┐   │
│      │ RemoteYansWifiChannelStub │  │ YansWifiChannel      │   │
│      │ -> MPI to Rank 0          │  │ (Local, Real Frames) │   │
│      │ -> High traffic (50Mbps)  │  │ -> PCAP Capture ✓    │   │
│      │ -> No PCAP ✗              │  │ -> Light (1-5Mbps)   │   │
│      └───────────────────────────┘  └──────────────────────┘   │
└─────────────────────────────────────────────────────────────────┘
\end{verbatim}
\caption{Hybrid dual-network architecture for MPI + PCAP}
\label{fig:hybrid-architecture}
\end{figure}

\paragraph{Network 1: MPI Performance Network}
\begin{itemize}
    \item Uses \texttt{RemoteYansWifiChannelStub}
    \item High data rates (50 Mbps for streaming, 30 Mbps for gaming, etc.)
    \item Distributed across machines via MPI
    \item Tests scalability and multi-machine performance
    \item IP subnet: 192.168.x.x (where x = rank ID)
\end{itemize}

\paragraph{Network 2: Local Monitoring Network}
\begin{itemize}
    \item Uses standard \texttt{YansWifiChannel}
    \item Light traffic (1-5 Mbps) for protocol monitoring
    \item Runs locally on each rank
    \item Generates PCAP files with real 802.11 frames
    \item IP subnet: 10.x.x.x (where x = rank ID)
\end{itemize}

\subsection{Implementation Details}

\subsubsection{Dual Interface Configuration}

Each simulated device gets two network interfaces:

\begin{lstlisting}[style=cpp, caption={Dual-network device creation}]
// Create MPI channel stub for performance testing
Ptr<RemoteYansWifiChannelStub> mpiChannelStub = 
    CreateObject<RemoteYansWifiChannelStub>();
mpiChannelStub->SetRemoteChannelRank(0);
mpiChannelStub->SetLocalDeviceRank(systemId);

// Create local channel for PCAP capture
YansWifiChannelHelper localChannelHelper = 
    YansWifiChannelHelper::Default();
localChannelHelper.AddPropagationLoss(
    "ns3::LogDistancePropagationLossModel",
    "Exponent", DoubleValue(3.0));
Ptr<YansWifiChannel> localChannel = 
    localChannelHelper.Create();

// Configure WiFi PHY for both networks
YansWifiPhyHelper mpiPhy, localPhy;
mpiPhy.SetChannel(mpiChannelStub);
localPhy.SetChannel(localChannel);

// Create devices on both networks
WifiHelper wifi;
wifi.SetStandard(WIFI_STANDARD_80211n);
WifiMacHelper mac;
mac.SetType("ns3::AdhocWifiMac");

NetDeviceContainer mpiDevices = 
    wifi.Install(mpiPhy, mac, nodes);
NetDeviceContainer localDevices = 
    wifi.Install(localPhy, mac, nodes);

// Enable PCAP only on local monitoring devices
localPhy.EnablePcapAll("home-wifi-rank" + 
                       std::to_string(systemId), true);
\end{lstlisting}

\subsubsection{IP Address Space Separation}

The two networks use distinct IP subnets to avoid routing conflicts:

\begin{lstlisting}[style=cpp, caption={Dual IP stack configuration}]
InternetStackHelper stack;
stack.Install(nodes);

// MPI network addressing (192.168.x.0/24)
Ipv4AddressHelper mpiAddress;
std::ostringstream mpiSubnet;
mpiSubnet << "192.168." << systemId << ".0";
mpiAddress.SetBase(mpiSubnet.str().c_str(), 
                   "255.255.255.0");
Ipv4InterfaceContainer mpiInterfaces = 
    mpiAddress.Assign(mpiDevices);

// Local monitoring network (10.x.0.0/24)
Ipv4AddressHelper localAddress;
std::ostringstream localSubnet;
localSubnet << "10." << systemId << ".0.0";
localAddress.SetBase(localSubnet.str().c_str(), 
                     "255.255.255.0");
Ipv4InterfaceContainer localInterfaces = 
    localAddress.Assign(localDevices);
\end{lstlisting}

\subsubsection{Traffic Distribution}

Applications send different traffic on each network:

\begin{lstlisting}[style=cpp, caption={Dual-network traffic generation}]
// Heavy MPI traffic for performance testing (50 Mbps)
OnOffHelper mpiTraffic("ns3::UdpSocketFactory",
    InetSocketAddress(mpiInterfaces.GetAddress(1), 8080));
mpiTraffic.SetAttribute("OnTime", 
    StringValue("ns3::ConstantRandomVariable[Constant=1]"));
mpiTraffic.SetAttribute("OffTime", 
    StringValue("ns3::ConstantRandomVariable[Constant=0]"));
mpiTraffic.SetAttribute("DataRate", DataRateValue(50000000));
mpiTraffic.SetAttribute("PacketSize", UintegerValue(1400));

ApplicationContainer mpiApp = mpiTraffic.Install(nodes.Get(0));
mpiApp.Start(Seconds(2.0));
mpiApp.Stop(Seconds(simTime - 1.0));

// Light monitoring traffic for PCAP (5 Mbps)
OnOffHelper localTraffic("ns3::UdpSocketFactory",
    InetSocketAddress(localInterfaces.GetAddress(1), 9090));
localTraffic.SetAttribute("DataRate", DataRateValue(5000000));
localTraffic.SetAttribute("PacketSize", UintegerValue(1200));

ApplicationContainer localApp = localTraffic.Install(nodes.Get(0));
localApp.Start(Seconds(2.0));
localApp.Stop(Seconds(simTime - 1.0));
\end{lstlisting}

\subsection{Traffic Patterns and Synchronization}

The dual networks operate independently but share the same simulation timeline:

\begin{itemize}
    \item \textbf{Same devices}: Each node participates in both networks
    \item \textbf{Same mobility}: Position updates affect both channels
    \item \textbf{Independent applications}: Different traffic generators
    \item \textbf{Synchronized time}: Both use the same \texttt{Simulator} clock
\end{itemize}

For example, in our home WiFi scenario:
\begin{itemize}
    \item \textbf{Smart TV}: 50 Mbps on MPI network (performance test) + 5 Mbps on local network (PCAP)
    \item \textbf{IoT devices}: 1 Mbps on MPI network + 1 Mbps on local network
\end{itemize}

\subsection{Benefits of Hybrid Approach}

\subsubsection{Performance + Observability}

Researchers get both capabilities:
\begin{itemize}
    \item \textbf{MPI network}: Tests system scalability with realistic traffic loads across machines
    \item \textbf{Local network}: Provides detailed protocol traces for analysis
\end{itemize}

\subsubsection{Debugging Distributed Scenarios}

When distributed simulations encounter problems, PCAP files reveal:
\begin{itemize}
    \item Whether packets are formed correctly (802.11 headers, frame types)
    \item MAC-level behavior (RTS/CTS, ACK, Block ACK)
    \item Timing issues (backoff, collision avoidance)
    \item Protocol state machines (authentication, association for infrastructure mode)
\end{itemize}

\subsubsection{Validation Against Standards}

Captured frames can be compared against IEEE 802.11 standards to validate:
\begin{itemize}
    \item Frame formats and field values
    \item Management frame sequences
    \item Error handling procedures
    \item Rate adaptation behavior
\end{itemize}

\subsubsection{Publication-Quality Traces}

PCAP files can be:
\begin{itemize}
    \item Opened in Wireshark for visual analysis
    \item Processed with \texttt{tcpdump} for statistics
    \item Converted to publication figures showing protocol exchanges
    \item Shared with reviewers to demonstrate correctness
\end{itemize}

\subsection{Performance Impact}

The hybrid approach adds overhead from:
\begin{enumerate}
    \item Additional network interface per device
    \item Duplicate MAC/PHY processing
    \item PCAP file I/O
\end{enumerate}

However, our measurements (Section~\ref{sec:results}) show this is acceptable:
\begin{itemize}
    \item Light monitoring traffic (1-5 Mbps) adds minimal computational load
    \item PCAP writes are buffered and asynchronous
    \item Memory increase is proportional to device count (not O(n²))
\end{itemize}

For our 8-device home WiFi scenario, the overhead is:
\begin{itemize}
    \item CPU: <10\% increase compared to MPI-only
    \item Memory: ~50 MB additional per rank for PCAP buffers
    \item Disk: 50-200 MB PCAP files per device (valuable for analysis)
\end{itemize}

\subsection{Alternative Considered}

\paragraph{Post-Processing Approach}
An alternative would be logging MPI messages and reconstructing PCAP files post-simulation. However, this requires:
\begin{itemize}
    \item Custom logging format
    \item Complete MPI message capture (storage overhead)
    \item Complex reconstruction logic
    \item Loss of real-time debugging capability
\end{itemize}

Our dual-network approach generates standard PCAP files in real-time, compatible with existing analysis tools, at acceptable computational cost.

\subsection{Applicability}

The hybrid approach is most beneficial when:
\begin{itemize}
    \item Validating new protocols or implementations
    \item Debugging unexpected distributed behavior
    \item Demonstrating correctness to reviewers
    \item Teaching/demonstrating WiFi protocols
\end{itemize}

For pure performance testing without protocol analysis, the MPI network alone suffices. The architecture allows easy switching between modes.
